% !TEX program = xelatex
% !TEX encoding = UTF-8

\documentclass[12pt, a4paper]{article}
\usepackage[UTF8]{ctex}
\usepackage{amsmath}
\usepackage{amssymb}
\usepackage{geometry}
\usepackage{enumitem}
\usepackage{fancyhdr}
\usepackage{titlesec}
\usepackage{xcolor}
\usepackage{listings}
\usepackage{framed}
\usepackage{tikz}

% 页面设置
\geometry{left=2.5cm, right=2.5cm, top=2.5cm, bottom=2.5cm}

% 页眉页脚
\pagestyle{fancy}
\fancyhf{}
\rhead{专升本数据结构讲义}
\lhead{章节精讲}
\cfoot{\thepage}

% 标题格式
\titleformat{\section}{\large\bfseries}{\thesection}{1em}{}
\titleformat{\subsection}{\normalsize\bfseries}{\thesubsection}{1em}{}

% 蓝色边框环境
\newenvironment{bluebox}{\begin{framed}\begin{minipage}{\linewidth}}{\end{minipage}\end{framed}}

% 代码样式
\lstset{
    backgroundcolor=\color{lightgray!20},
    basicstyle=\ttfamily\small,
    numbers=left,
    numberstyle=\tiny,
    frame=single,
    breaklines=true
}

\title{\textbf{第三章 栈和队列}}
\author{}
\date{}

\begin{document}

\maketitle

\section*{本章知识脉络}

栈和队列是两种重要的线性数据结构，它们都是操作受限的线性表。栈是一种后进先出（LIFO，Last In First Out）的数据结构，只允许在表的一端进行插入和删除操作。队列是一种先进先出（FIFO，First In First Out）的数据结构，允许在表的一端插入、在另一端删除。这两种数据结构在实际计算机程序中有广泛的应用，是算法设计中的重要工具。

栈和队列虽然操作受限，但这种限制带来了独特的功能特性。栈的LIFO特性使得它特别适合用于需要“记忆”功能的场景，如函数调用、表达式求值、括号匹配等。队列的FIFO特性使得它适合用于需要“先来先服务”的场景，如任务调度、打印队列、广度优先搜索等。

本章的重点内容包括：栈和队列的定义与基本操作、顺序栈和链栈的实现、顺序队列和循环队列的实现、栈和队列的应用、栈和队列在算法中的应用。难点包括：循环队列的判满判空条件、栈的出栈序列合法性判断、栈和队列的应用场景分析。

\section{栈的定义与特性}

栈（Stack）是限定仅在表尾进行插入和删除操作的线性表。栈的插入操作称为进栈（Push）或入栈，删除操作称为出栈（Pop）或退栈。允许进行操作的一端称为栈顶（Top），另一端称为栈底（Bottom）。不含任何元素的栈称为空栈。

栈的抽象数据模型可以描述为：数据对象是n个同类型数据元素的有限序列；数据关系是除了栈顶元素外，每个元素有且仅有一个直接前驱；栈顶元素是最后进入栈的元素，最先被删除。栈的基本操作包括：初始化栈、判断栈是否为空、获取栈顶元素、入栈、出栈、销毁栈等。

栈的最大的特点是后进先出（LIFO）。想象一下只有一个口的羽毛球筒：最先放入的球在筒底，最后放入的球在筒口。取球时，最先取出的只能是最后放入的球。这就是栈的LIFO特性的直观体现。这个特性使得栈成为一种非常有用的数据结构，在计算机科学的诸多领域都有广泛应用。

栈是一种线性结构，但与普通线性表不同的是，栈是一种操作受限的线性表。普通线性表可以在任意位置插入和删除，而栈只能在栈顶进行插入和删除操作。这种限制虽然降低了灵活性，但使得栈的实现更加简单高效，同时也带来了独特的功能特性。

栈的存储结构可以分为顺序存储和链式存储两种。顺序存储的栈称为顺序栈，利用一组连续的存储单元存放栈中的元素。链式存储的栈称为链栈，利用链表实现。两种实现方式各有优缺点：顺序栈的访问效率高，但需要预先分配固定大小的空间；链栈的空间可以动态扩展，但每个结点需要额外的指针空间。

\section{栈的顺序存储实现}

顺序栈是使用数组实现的栈结构。通常用一个一维数组和一个指向栈顶的变量来实现。栈顶指针可以是数组下标（通常用top表示，top指向当前栈顶元素的位置），也可以是索引（top指向下一个可入栈的位置）。不同的实现方式在判空判满条件上有所不同。

顺序栈的存储结构定义通常包括：一个存放数据元素的数组data[MAXSIZE]（MAXSIZE是栈的最大容量）；一个表示栈顶位置的变量top。栈顶指针的初始值通常为-1或0，取决于具体的实现方式。当top为-1时，表示栈为空；当top为MAXSIZE-1时，表示栈已满。

顺序栈的入栈操作（Push）步骤：首先判断栈是否已满，如果已满则返回错误或溢出；然后将栈顶指针加1；最后将新元素存入栈顶位置。入栈操作的时间复杂度为O(1)。

顺序栈的出栈操作（Pop）步骤：首先判断栈是否为空，如果为空则返回错误；然后取出栈顶元素；最后将栈顶指针减1。出栈操作的时间复杂度为O(1)。

顺序栈的优点是实现简单、访问效率高，缺点是需要预先分配固定大小的存储空间，可能造成空间浪费或不足。在实际应用中，可以采用动态扩容的方式来解决空间不足的问题：当栈满时，分配一个更大的数组，将原数组的元素复制到新数组中，然后释放原数组。

\section{栈的链式存储实现}

链栈是使用链表实现的栈结构。链栈的栈顶指针指向链表的头结点（带头结点的链栈）或第一个结点（不带头结点的链栈）。由于栈的插入和删除操作都在栈顶进行，链表的头部就是栈顶，因此链栈的入栈和出栈操作都非常高效。

链栈的结点结构包括：数据域（存放数据元素）和指针域（指向下一个结点）。链栈的栈顶指针指向栈的第一个结点（对于不带头结点的实现）或指向头结点（对于带头结点的实现）。

链栈的入栈操作（Push）步骤：创建新结点；将新结点的指针指向当前栈顶结点；将栈顶指针指向新结点。入栈操作的时间复杂度为O(1)。

链栈的出栈操作（Pop）步骤：判断栈是否为空；取出栈顶元素；将栈顶指针指向栈顶元素的下一个结点；释放原栈顶结点的空间。出栈操作的时间复杂度为O(1)。

链栈的优点是空间可以动态扩展，不会出现上溢出的问题；缺点是每个结点需要额外的指针空间，存储密度较低。由于栈操作只在栈顶进行，链栈的链表只需要单链表即可，不需要双向链表或循环链表。

链栈不需要头结点，因为栈的操作只在链表头部进行。带头结点的链栈反而增加了复杂性，因为头结点不存储实际数据，而且入栈出栈操作都需要跳过头结点。因此，链栈通常使用不带头结点的单链表实现。

栈的应用非常广泛。在计算机科学中，栈主要用于以下几个方面：函数调用和递归调用时的函数返回地址和局部变量的存储；表达式求值和语法分析；括号匹配检测；深度优先搜索；函数值的逆序输出等。

\section{队列的定义与特性}

队列（Queue）是限定仅在表的一端进行插入操作、在另一端进行删除操作的线性表。队列的插入操作称为入队（Enqueue）或进队，在队尾进行；队列的删除操作称为出队（Dequeue）或退队，在队首进行。允许插入的一端称为队尾（Rear），允许删除的一端称为队首（Front）。

队列的抽象数据模型可以描述为：数据对象是n个同类型数据元素的有限序列；数据关系是除了队首元素外，每个元素有且仅有一个直接前驱；队首元素是最先进入队列的元素，最先被删除。队列的基本操作包括：初始化队列、判断队列是否为空、获取队首元素、入队、出队、销毁队列等。

队列的最大特点是先进先出（FIFO）。想象一下排队买票：最先排队的人最先买到票，后来的人后买到票。这就是队列的FIFO特性的直观体现。这个特性使得队列特别适合用于需要“先来先服务”的场景。

队列是一种操作受限的线性表，与栈类似，都是在线性表的基础上增加了操作限制。队列的基本操作只有两个：入队和出队。这种限制使得队列的实现更加简单高效，也带来了独特的功能特性。

队列的存储结构也可以分为顺序存储和链式存储两种。顺序存储的队列称为顺序队列（或循环队列），链式存储的队列称为链式队列。循环队列是顺序队列的一种重要形式，它巧妙地解决了普通顺序队列的“假溢出”问题。

队列在计算机科学中有广泛的应用。在操作系统中，进程调度使用队列管理就绪进程；在打印机管理中，打印任务按队列顺序执行；在广度优先搜索算法中，队列用于存储待访问的结点；在消息队列中，队列用于实现进程间的通信等。

\section{顺序队列与循环队列}

顺序队列是使用数组实现的队列结构。通常用一个一维数组和两个指针（队首指针front和队尾指针rear）来实现。队首指针指向队首元素的前一个位置，队尾指针指向队尾元素的位置（也有其他实现方式）。

顺序队列的入队操作步骤：将新元素存入rear指向的位置；将rear加1。出队操作步骤：取出front+1位置的元素；将front加1。但是，这种实现方式存在“假溢出”问题：当队列未满时，rear可能已经到达数组末尾，无法继续入队。

循环队列是解决顺序队列“假溢出”问题的有效方法。循环队列将数组的首尾相连，形成一个环。当rear到达数组末尾时，下一步移动到数组开头；同样，当front到达数组末尾时，下一步也移动到数组开头。

循环队列的判空和判满条件是一个重要的考点。由于循环队列中front和rear相等时既可能表示队列为空，也可能表示队列为满（假设队列满时rear与front也相等），因此需要采用特殊的方法来区分这两种情况。常用的方法有三种：牺牲一个存储单元法（常用）、设置计数器法、设置标志位法。

牺牲一个存储单元法是最常用的方法。具体做法是：约定队列满的条件是(rear+1) \% MAXSIZE == front，队列空的条件是front == rear。这种方法牺牲了一个存储单元，但实现简单。在MAXSIZE个存储单元的数组中，最多可以存储MAXSIZE-1个元素。

循环队列的入队操作步骤：判断队列是否已满；如果未满，将新元素存入rear位置；将rear循环加1（rear = (rear + 1) \% MAXSIZE）。出队操作步骤：判断队列是否为空；如果不为空，取出front+1位置的元素；将front循环加1（front = (front + 1) \% MAXSIZE）。

循环队列的优点是解决了顺序队列的假溢出问题，空间利用率高；缺点是需要预先分配固定大小的存储空间。循环队列是考试的重点内容，考生需要熟练掌握循环队列的判空判满条件以及入队出队操作的实现。

\section{链式队列}

链式队列是使用链表实现的队列结构。链式队列需要两个指针：队首指针front指向队首结点，队尾指针rear指向队尾结点。为了操作方便，链式队列通常带头结点：头结点不存储实际数据，仅作为队列的入口。

链式队列的结点结构包括：数据域（存放数据元素）和指针域（指向下一个结点）。队首结点的数据是队首元素，队尾结点的数据是队尾元素。

链式队列的入队操作（Enqueue）步骤：创建新结点；将新结点的指针置为NULL；将rear结点的指针指向新结点；更新rear指针指向新结点。入队操作的时间复杂度为O(1)。

链式队列的出队操作（Dequeue）步骤：判断队列是否为空；如果不为空，取出front->next结点的数据；将front指针向后移动一个结点；如果出队后队列为空，需要将rear指针也置空。出队操作的时间复杂度为O(1)。

链式队列的优点是空间可以动态扩展，不会出现上溢出的问题；每个结点需要额外的指针空间。链式队列不会出现“假溢出”问题，因为链表可以动态增长。

链式队列的判空条件是front == rear（带头结点的情况，此时front和rear都指向头结点）。链式队列不需要像循环队列那样考虑判空判满的问题，因为链表可以动态扩展。

\section{栈与队列的对比}

栈和队列都是操作受限的线性表，但它们的操作特性和应用场景有显著区别。理解这些区别对于正确选择和使用这两种数据结构至关重要。

从操作特性来看，栈是后进先出（LIFO），最后入栈的元素最先出栈；队列是先进先出（FIFO），最先入队的元素最先出队。这个根本区别决定了两者不同的应用场景。

从操作限制来看，栈只能在栈顶进行插入和删除操作，是单端操作；队列在一端插入、在另一端删除，是双端操作但方向固定。

从存储实现来看，栈和队列都可以用顺序存储或链式存储实现。顺序栈和顺序队列（循环队列）的访问效率高，但需要预先分配空间；链栈和链式队列的空间可以动态扩展。

从时间复杂度来看，栈和队列的入栈、入队、出栈、出队操作的时间复杂度都是O(1)。这是因为栈和队列都是操作受限的数据结构，插入和删除操作都只涉及端点位置的修改，不需要移动元素。

从空间复杂度来看，顺序存储的空间复杂度为O(n)，链式存储的空间复杂度也为O(n)，但链式存储每个结点需要额外的指针空间。

栈和队列的选择取决于具体的应用场景。如果需要实现“后进先出”的功能，选择栈；如果需要实现“先进先出”的功能，选择队列。栈适合用于函数调用、表达式求值、括号匹配等场景；队列适合用于任务调度、广度优先搜索、消息队列等场景。

\begin{bluebox}{本章核心考点总结}

	\textbf{【核心考点一：栈的LIFO特性】}

	栈的核心特性是后进先出（Last In First Out）。栈的插入操作在栈顶进行，删除操作也在栈顶进行。最先入栈的元素在栈底，最后入栈的元素在栈顶，最先被删除。这个特性是理解栈和应用栈的基础。

	\textbf{【核心考点二：队列的FIFO特性】}

	队列的核心特性是先进先出（First In First Out）。队列的插入操作在队尾进行，删除操作在队首进行。最先入队的元素最先出队，最后入队的元素最后出队。这个特性使得队列适合用于“先来先服务”的场景。

	\textbf{【核心考点三：循环队列的判空判满条件】}

	这是本章最重要的考点。常用方法是牺牲一个存储单元法：队列满的条件是(rear+1) \% MAXSIZE == front，队列空的条件是front == rear。考生需要熟练掌握这个判定方法，并能够根据具体实现方式推导判空判满条件。

	\textbf{【核心考点四：栈和队列的操作复杂度】}

	栈和队列的所有基本操作（入栈、出栈、入队、出队、获取栈顶/队首元素、判断空/满）的时间复杂度都是O(1)。这是栈和队列的重要特点，也是它们被广泛应用的原因之一。

	\textbf{【核心考点五：栈和队列的应用】}

	栈的主要应用包括：函数调用（递归）、表达式求值、括号匹配、深度优先搜索等。队列的主要应用包括：任务调度、广度优先搜索、打印队列、消息队列等。理解这些应用有助于加深对栈和队列特性的理解。
\end{bluebox}

\section{典型例题精讲}

\subsection{选择题精析}

\begin{bluebox}{答案与深度解析}
	\textbf{答案：B}

	\textbf{【核心认知框架】}

	本题考察循环队列的判满条件。循环队列是顺序队列的改进形式，通过将数组首尾相连解决了“假溢出”问题。但循环队列的判空和判满需要特殊处理，这是数据结构课程的重点和难点。

	\textbf{【选项原理深度拆解】}

	\item \textbf{循环队列的基本原理}
	\begin{itemize}
		\item 循环队列将顺序队列的数组首尾相连，形成逻辑上的环形结构
		\item 当rear到达数组末尾时，下一步移动到数组开头（rear = 0）
		\item 通过取模运算实现循环：(rear + 1) \% MAXSIZE
	\end{itemize}

	\item \textbf{判空判满的特殊处理}
	\begin{itemize}
		\item 如果简单地用front == rear表示空、rear == MAXSIZE表示满
		\item 当队列满时，rear == MAXSIZE，此时front == rear == MAXSIZE
		\item 这导致无法区分队列空和队列满两种状态
		\item 需要采用特殊方法解决这个问题
	\end{itemize}

	\item \textbf{牺牲一个存储单元法}
	\begin{itemize}
		\item 约定：队列满的条件是(rear + 1) \% MAXSIZE == front
		\item 队列空的条件是front == rear
		\item 这种方法牺牲了一个存储单元，最多存储MAXSIZE-1个元素
		\item 这是最常用的方法
	\end{itemize}

	\textbf{【答案解析】}

	根据牺牲一个存储单元法，循环队列满的条件是(rear + 1) \% MAXSIZE == front。题目中MAXSIZE为6，rear为4，代入计算：(4+1) \% 6 = 5，不等于front(1)，因此队列未满。答案为B。

	\textbf{【应背诵知识点】}

	循环队列判满：(rear + 1) \% MAXSIZE == front；判空：front == rear。
\end{bluebox}

\begin{bluebox}{答案与深度解析}
	\textbf{答案：A}

	\textbf{【核心认知框架】}

	本题考察栈的输出序列合法性问题。对于一个输入序列，经过栈的输出可能产生多种不同的输出序列。理解栈的LIFO特性是解决此类问题的关键。

	\textbf{【选项原理深度拆解】}

	\item \textbf{栈的LIFO特性}
	\begin{itemize}
		\item 入栈顺序：1, 2, 3, 4
		\item 出栈时，最后入栈的4最先出来
		\item 3必须在4之后出来（因为4在栈顶）
		\item 2必须在3之后出来
		\item 1必须在2之后出来
	\end{itemize}

	\item \textbf{验证选项A：4, 1, 3, 2}
	\begin{itemize}
		\item 第一个输出是4：入栈1,2,3,4，出栈4 ✓
		\item 第二个输出是1：栈中有3,2,1，但1在栈底，需要先出3和2
		\item 但题目输出1在3和2之前，这是不可能的 ✗
		\item 因此A不合法
	\end{itemize}

	\item \textbf{验证选项B：4, 3, 2, 1}
	\begin{itemize}
		\item 入栈1,2,3,4
		      \出栈4,3,2,1，完全符合LIFO ✓
	\end{itemize}

	\item \textbf{验证选项C：4, 3, 1, 2}
	\begin{itemize}
		\item 入栈1,2,3,4，出栈4,3
		\item 栈中剩余2,1
		\item 下一个输出1（栈顶是2，但输出1需要2先出）✗
	\end{itemize}

	\item \textbf{验证选项D：4, 2, 1, 3}
	\begin{itemize}
		\item 入栈1,2,3,4，出栈4,2
		\item 栈中剩余3,1
		\item 下一个输出1，再输出3 ✓
	\end{itemize}

	\textbf{【答案解析】}

	选项A中，4先出栈后，栈中还剩3,2,1，此时1在栈底，要输出1必须先输出2和3，但输出序列中1在3和2之前，因此不可能。

	\textbf{【快速判断方法】}

	对于栈的输出序列，可以用一个简单的原则判断：在输出序列中，如果某个元素后面出现了比它更小的元素，那么这些更小的元素在该元素出栈时必须都在栈中。例如，4后面是1,3,2，当4出栈时，1,3,2都在栈中，1在栈底，要输出1必须先输出3和2。
\end{bluebox}

\begin{bluebox}{答案与深度解析}
	\textbf{答案：C}

	\textbf{【核心认知框架】}

	本题考察循环队列的原理和假溢出问题。循环队列通过将队列的首尾相连形成环状结构，解决了普通顺序队列的假溢出问题。理解循环队列的原理是解题的关键。

	\textbf{【选项原理深度拆解】}

	\item \textbf{普通顺序队列的问题}
	\begin{itemize}
		\item 随着入队和出队的进行，front和rear不断增大
		\item 当rear达到数组末尾时，即使队列实际未满，也无法继续入队
		      \frontend这就是“假溢出”问题
	\end{itemize}

	\item \textbf{循环队列的解决方案}
	\begin{itemize}
		\item 将数组首尾相连，形成逻辑上的环
		\item 通过取模运算实现循环：index = (index + 1) \% MAXSIZE
		\item front和rear可以循环移动
	\end{itemize}

	\item \textbf{循环队列的空间利用}
	\begin{itemize}
		\item 数组大小为MAXSIZE时，最多存储MAXSIZE个元素
		\item 但实际可用空间为MAXSIZE-1个（牺牲一个单元判满）
	\end{itemize}

	\textbf{【答案解析】}

	循环队列解决了普通顺序队列的假溢出问题，使得队列可以循环使用。答案为C。

	\textbf{【应背诵知识点】}

	循环队列的本质是用取模运算实现下标的循环，使得front和rear可以回到数组开头继续使用。
\end{bluebox}

\begin{bluebox}{答案与深度解析}
	\textbf{答案：D}

	\textbf{【核心认知框架】}

	本题考察栈和队列的特性对比。栈是后进先出（LIFO），队列是先进先出（FIFO）。这个根本区别决定了两者完全不同的操作特性和应用场景。

	\textbf{【选项原理深度拆解】}

	\item \textbf{栈的特性}
	\begin{itemize}
		\item 后进先出（Last In First Out，LIFO）
		\item 插入和删除都在栈顶进行
		\item 适用于需要“记忆”功能的场景
	\end{itemize}

	\item \textbf{队列的特性}
	\begin{itemize}
		\item 先进先出（First In First Out，FIFO）
		\item 插入在队尾，删除在队首
		\item 适用于“先来先服务”场景
	\end{itemize}

	\item \textbf{选项分析}
	\begin{itemize}
		\item[A.] 栈和队列都是线性表——正确，两者都是线性结构
		\item[B.] 栈是后进先出——正确
		\item[C.] 队列是先进先出——正确
		\item[D.] 栈和队列的插入删除操作位置相同——错误，栈在一端，队列在两端
	\end{itemize}

	\textbf{【答案解析】}

	栈的插入和删除都在栈顶进行，队列的插入在队尾、删除在队首，两者的操作位置不同。答案为D。
\end{bluebox}

\begin{bluebox}{答案与深度解析}
	\textbf{答案：A}

	\textbf{【核心认知框架】}

	本题考察队列操作的特点。队列的删除操作只能在队首进行，这是队列FIFO特性的体现。理解队列的基本操作是解题的关键。

	\textbf{【选项原理深度拆解】}

	\item \textbf{队列的基本操作}
	\begin{itemize}
		\item 入队（Enqueue）：在队尾插入
		\item 出队（Dequeue）：在队首删除
		\item 只能在一端插入，另一端删除
	\end{itemize}

	\item \textbf{选项分析}
	\begin{itemize}
		\item[A.] 插入和删除位置限制——正确
		\item[B.] 插入和删除时间复杂度——均为O(1)
		\item[C.] 插入和删除位置——队首和队尾，不是都有限制
		\item[D.] 插入删除位置限制——队尾插入，队首删除
	\end{itemize}

	\textbf{【答案解析】}

	队列的插入只能在队尾进行，删除只能在队首进行，这是队列的基本特性。答案为A。
\end{bluebox}

\subsection{填空题精析}

\begin{enumerate}
	\item \textbf{栈是（后进先出）或（先进后出）的线性表。}

	      解析：栈最重要的特性就是后进先出（Last In First Out，LIFO），也称为先进后出。最先入栈的元素最后出栈，最后入栈的元素最先出栈。

	\item \textbf{队列是（先进先出）的线性表。}

	      解析：队列最重要的特性是先进先出（First In First Out，FIFO）。最先入队的元素最先出队，最后入队的元素最后出队。

	\item \textbf{循环队列的队首指针指向第一个（元素）的前一个位置，队尾指针指向（最后一个元素）的位置。}

	      解析：这是循环队列的一种常见实现方式。队首指针指向队首元素的前一个位置，队尾指针指向队尾元素的位置。

	\item \textbf{设循环队列的容量为MAXSIZE，队首指针为front，队尾指针为rear，则队空的条件是（front == rear），队满的条件是（(rear+1) \% MAXSIZE == front）。}

	      解析：这是牺牲一个存储单元法的判空判满条件。队空时front == rear，队满时rear的下一个位置是front。

	\item \textbf{在具有n个单元的循环队列中，最多能存放（n-1）个元素。}

	      解析：由于采用牺牲一个存储单元法来区分队空和队满，最多只能存储n-1个元素。

	\item \textbf{栈的插入操作称为（进栈或入栈），删除操作称为（出栈或退栈）。}

	      解析：栈的基本操作名称需要掌握。

	\item \textbf{队列的插入操作称为（入队或进队），删除操作称为（出队或退队）。}

	      解析：队列的基本操作名称需要掌握。
\end{enumerate}

\subsection{判断题精析}

\begin{enumerate}
	\item \textbf{栈的操作原则是先进后出。（正确）}

	      解析：栈是后进先出（LIFO）的数据结构。正确表述是“后进先出”或“先进后出”。

	\item \textbf{队列的操作原则是先进先出。（正确）}

	      解析：队列是先进先出（FIFO）的数据结构。最先进入队列的元素最先离开队列。

	\item \textbf{循环队列是为了解决顺序队列的假溢出问题。（正确）}

	      解析：循环队列通过将数组首尾相连，解决了顺序队列的“假溢出”问题。这是循环队列产生的原因。

	\item \textbf{栈和队列都是线性结构。（正确）}

	      解析：栈和队列都是操作受限的线性表，本质上都是线性结构。它们是在线性表基础上增加了操作限制。

	\item \textbf{栈和队列都可以用顺序存储结构实现。（正确）}

	      解析：栈可以用顺序栈实现，队列可以用顺序队列（循环队列）实现。两种数据结构都可以用顺序存储实现。

	\item \textbf{栈和队列都可以用链式存储结构实现。（正确）}

	      解析：栈可以用链栈实现，队列可以用链式队列实现。链式存储使得存储空间可以动态扩展。

	\item \textbf{循环队列的队空和队满的判定条件相同。（错误）}

	      解析：循环队列的队空和队满判定条件不同。常用方法是牺牲一个存储单元法：队空是front == rear，队满是(rear+1) \% MAXSIZE == front。
\end{enumerate}

\section{本章知识体系图}

\begin{center}
	\begin{tikzpicture}[scale=0.85, every node/.style={scale=0.85}]
		\tikzstyle{bx} = [rectangle, rounded corners, minimum width=2.5cm, minimum height=0.8cm, text centered, draw=black, fill=blue!20]
		\tikzstyle{ar} = [->, >=stealth, thick]

		\node[bx] (sq) at (-4,0) {栈（Stack）};
		\node[bx] (queue) at (4,0) {队列（Queue）};
		\node[bx] (lif) at (-4,-1.5) {LIFO};
		\node[bx] (fif) at (4,-1.5) {FIFO};
		\node[bx] (push) at (-6,-3) {进栈};
		\node[bx] (pop) at (-4,-3) {出栈};
		\node[bx] (top) at (-2,-3) {栈顶};
		\node[bx] (enq) at (2,-3) {入队};
		\node[bx] (deq) at (4,-3) {出队};
		\node[bx] (front) at (6,-3) {队首};
		\node[bx] (rear) at (8,-3) {队尾};

		\node[bx] (sseq) at (-5,-5) {顺序栈};
		\node[bx] (slink) at (-3,-5) {链栈};
		\node[bx] (qseq) at (3,-5) {循环队列};
		\node[bx] (qlink) at (5,-5) {链式队列};

		\draw[ar] (sq) -- (lif);
		\draw[ar] (queue) -- (fif);
		\draw[ar] (lif) -- (push);
		\draw[ar] (lif) -- (pop);
		\draw[ar] (lif) -- (top);
		\draw[ar] (fif) -- (enq);
		\draw[ar] (fif) -- (deq);
		\draw[ar] (fif) -- (front);
		\draw[ar] (fif) -- (rear);
		\draw[ar] (sq) -- (sseq);
		\draw[ar] (sq) -- (slink);
		\draw[ar] (queue) -- (qseq);
		\draw[ar] (queue) -- (qlink);
	\end{tikzpicture}
\end{center}

\section{本章学习建议}

本章内容相对简单，但概念较多，需要准确记忆。建议考生从以下方面学习：

第一，准确理解栈和队列的定义。栈是后进先出（LIFO），队列是先进先出（FIFO）。这个核心特性决定了它们的应用场景。

第二，重点掌握循环队列的判空判满条件。这是本章的难点和重点，需要通过练习熟练掌握。牺牲一个存储单元法是最常用的方法：队满(rear+1)\%MAXSIZE == front，队空front == rear。

第三，理解栈和队列的应用场景。栈适用于函数调用、表达式求值、括号匹配等；队列适用于任务调度、广度优先搜索等。理解应用场景有助于加深对特性的理解。

第四，多做练习，熟能生巧。栈的输出序列合法性判断是常见的题型，需要多做练习掌握判断方法。

第五，注意区分易混淆概念。栈的“先进后出”和队列的“先进先出”不要混淆；循环队列的判空判满条件要准确记忆。

\end{document}